\section{Introduction}
Many predictors output a \emph{composition}: a non-negative vector summing to one, such as softmax probabilities, topic proportions, or the empirical label distribution of a panel of annotators. A user of such a predictor needs to know how far the true composition may be from the prediction. Split conformal prediction \citep{vovk2005algorithmic,papadopoulos2002inductive,lei2018distribution} offers a distribution-free answer: given a non-conformity score $s(y,p)$ and $n$ exchangeable calibration pairs, the set $\{y: s(y,\hat p(x))\le \hat q\}$ covers the true target with probability at least $1-\alpha$. For classification, scores and sets for the \emph{label} are well studied \citep{sadinle2019least,romano2020classification,angelopoulos2021uncertainty,ding2023class}, and set-valued prediction of probability distributions has been proposed \citep{javanmardi2024conformalized}. What is missing is a clear, reproducible account of \emph{how to evaluate} conformal methods when the target itself lies on the simplex: which score, which calibration scheme, and which metric?

Validity is guaranteed by the recipe for \emph{every} score, so marginal coverage cannot discriminate between methods; the interesting differences are (i) the \emph{size} of the sets, which depends on the geometry chosen for the simplex \citep{aitchison1982statistical}, and (ii) \emph{conditional} behaviour, where no distribution-free guarantee is possible in general \citep{vovk2012conditional,barber2021limits}. Evaluating (i) and (ii) raises specific methodological problems: set ``size'' has no canonical measure on the simplex; worst-slice coverage metrics are biased downward even for perfectly valid methods; and the severity of conditional failure should be connected to a property of the data-generating process.

\paragraph{Research questions.} We study split conformal prediction with simplex-valued targets and ask:
\begin{itemize}[leftmargin=1.2em,itemsep=0pt,topsep=2pt]
\item \textbf{RQ1 (validity).} Does the finite-sample coverage law of split conformal prediction hold, quantitatively, for simplex targets and for scores of very different geometry (including ones that are discontinuous or unbounded at the boundary)?
\item \textbf{RQ2 (efficiency).} Does the choice among six natural scores ($\ell_2$, total variation, $\ell_\infty$, Hellinger, Aitchison, KL) change the size of the prediction set in a way that is stable across the dimension $K$ and concentration of the target distribution? And does the ranking depend on how size is measured?
\item \textbf{RQ3 (conditional behaviour).} When does plain split calibration give misleading conditional coverage, and which of two standard remedies (score normalisation by an estimated scale, Mondrian calibration on predicted-confidence bins) helps under which failure mode?
\item \textbf{RQ4 (robustness to shift).} On a real compositional target, what happens to the coverage guarantee under covariate shift between calibration and test, with and without weighting?
\end{itemize}
\paragraph{Contributions.} (1) A controlled synthetic benchmark in which the target is $\mathrm{Dirichlet}(\kappa(x)m(x))$, so that concentration, heteroscedasticity and mean mis-specification are separately controlled, with 2,400 independent runs of 18 conformal variants. (2) An exact reference for evaluating conformal methods: the Beta law of realised coverage, which we verify to Monte-Carlo precision for three scores, plus a \emph{null reference} for worst-slice coverage metrics, which we show is essential for reading them. (3) A closed-form variance identity for the Dirichlet that explains \emph{why} $\ell_2$-type scores are heteroscedastic on the simplex and predicts the failure of plain calibration; the experiment shows which remedy addresses which cause. (4) A real-data evaluation on CIFAR-10 with human soft labels (CIFAR-10H), including shift. We report negative and mixed findings explicitly: no score dominates across settings, and neither remedy dominates.

\section{Related Work}
\paragraph{Conformal prediction.} Our baseline is split (inductive) conformal prediction \citep{papadopoulos2002inductive,lei2018distribution}; the full treatment is in \citet{vovk2005algorithmic} and the tutorial of \citet{angelopoulos2021gentle}. The exact distribution of the realised coverage of split conformal given the calibration set is a Beta law \citep{vovk2012conditional}; we use it as a reference for evaluation, not as a new result. Distribution-free \emph{conditional} coverage is impossible in general \citep{barber2021limits}; practical surrogates are Mondrian (group-conditional) calibration \citep{vovk2005algorithmic,ding2023class}, locally-weighted scores \citep{lei2018distribution}, and quantile-based scores \citep{romano2019conformalized}; multivalid methods give group-conditional guarantees over richer families \citep{jung2022batch}. Covariate shift is handled by weighted conformal prediction when the likelihood ratio is known \citep{tibshirani2019covariate}. Jackknife+ \citep{barber2021jackknife} and conformal risk control \citep{angelopoulos2022conformal} are not evaluated here.
\paragraph{Conformal prediction for classifiers.} Least-ambiguous sets, APS, RAPS and clustered conformal prediction construct sets of \emph{labels} \citep{sadinle2019least,romano2020classification,angelopoulos2021uncertainty,ding2023class}. The target there is a discrete label; ours is a point in the simplex, so a prediction set is a region of the simplex. \citet{javanmardi2024conformalized} construct conformal credal sets (sets of distributions) for classification; \citet{cresswell2024disparate} study the downstream effect of label sets on human decision-making and group disparities. We do not reproduce those methods; our contribution is an evaluation study of generic split conformal sets on the simplex with several scores.
\paragraph{Compositional data.} Aitchison's log-ratio geometry \citep{aitchison1982statistical} motivates the Aitchison distance score we include; since it is undefined on the boundary we use an additive pseudo-count (Sec.~\ref{sec:method}) and examine its sensitivity on real data. CIFAR-10H \citep{peterson2019human} supplies human label distributions for the CIFAR-10 \citep{krizhevsky2009learning} test images.

\section{Problem Setup, Scores and Theory}\label{sec:method}
\paragraph{Setting.} Let $\Delta^{K-1}=\{y\in\mathbb R^K_{\ge0}:\sum_k y_k=1\}$. A predictor outputs $\hat p(x)\in\Delta^{K-1}$ for a target $y\in\Delta^{K-1}$. With a score $s:\Delta^{K-1}\times\Delta^{K-1}\to[0,\infty)$, calibration scores $S_i=s(y_i,\hat p(x_i))$, $i\le n$, and $\ell=\lfloor(n+1)\alpha\rfloor$, the split conformal threshold $\hat q$ is the $(n+1-\ell)$-th smallest $S_i$ and the set is $C(x)=\{y:s(y,\hat p(x))\le\hat q\}$.
\paragraph{Scores.} We use $\ell_2$: $\|y-p\|_2$; TV: $\tfrac12\|y-p\|_1$; $\ell_\infty$: $\|y-p\|_\infty$; Hellinger: $\|\sqrt y-\sqrt p\|_2/\sqrt2$; Aitchison: $\|\mathrm{clr}(y_\delta)-\mathrm{clr}(p_\delta)\|_2$ with $y_\delta=(y+\delta)/(1+K\delta)$ ($\delta=10^{-2}$ by default); KL: $\mathrm{KL}(y\|p)$. For $K=3$, TV and $\ell_\infty$ coincide exactly (the signed differences sum to zero), a fact visible as identical results in our experiments. Aitchison and KL are unbounded when $\hat p$ has near-zero entries where $y$ does not.
\paragraph{Calibration variants.} \emph{plain}: as above. \emph{normalised}: calibrate $S_i/\hat\sigma(x_i)$ where $\log\hat\sigma(x)$ is a ridge regression on features of $x$ fitted on the (disjoint) training split. \emph{Mondrian}: a separate threshold within each of $5$ bins of the predicted maximum probability (bin edges from the training split). \emph{Weighted} (real-data shift only): the weighted quantile of \citet{tibshirani2019covariate} with the known likelihood ratio.

\begin{proposition}[Exact coverage law; \citealp{vovk2012conditional}]\label{prop:beta}
If the calibration and test scores are exchangeable with a continuous distribution, then $C=\mathbb P(S_{n+1}\le\hat q\mid S_{1:n})\sim\mathrm{Beta}(n+1-\ell,\ell)$, with $\mathbb E C=\frac{n+1-\ell}{n+1}\ge1-\alpha$ and $\mathrm{Var}\,C=\frac{(n+1-\ell)\ell}{(n+1)^2(n+2)}$. Mondrian calibration gives the same law with $n$ replaced by the group size $n_g$, group by group.
\end{proposition}
\emph{Proof sketch.} $C=U_{(n+1-\ell)}$-distributed as the $(n+1-\ell)$-th order statistic of $n$ i.i.d.\ uniforms (probability integral transform), which is $\mathrm{Beta}(n+1-\ell,\ell)$. \qed\\
Consequences for \emph{evaluation}: (a) the sampling SD of realised coverage with $n_{\rm cal}=500$, $n_{\rm test}=1000$ is $\approx0.016$ (Beta part $0.0133$, binomial part $0.0095$), so single-run coverage deviations of $\pm 0.03$ are noise; (b) the \emph{expected} coverage is $0.9002$, slightly above $0.9$, and a method-comparison must not interpret differences of that size as meaningful.

\begin{proposition}[Worst-slice coverage is biased under perfect validity]\label{prop:null}
Let slices $A_1,\dots,A_B$ have $m$ test points each and let the conditional coverage of every slice equal the realised marginal coverage $C$ (perfect conditional validity for the given threshold). Then $\mathbb E\min_b \hat c_b<\mathbb E C$ strictly whenever $m<\infty$, where $\hat c_b\sim\mathrm{Bin}(m,C)/m$.
\end{proposition}
This is immediate (the minimum of $B\ge2$ non-degenerate random variables has strictly smaller mean). Its practical content is quantitative: with $B=5$ slices of $m=200$ points and $n=500$ the null value of the worst-slice coverage is $0.875$ for $1-\alpha=0.9$ (simulated from Prop.~\ref{prop:beta}, 20{,}000 replications), so \emph{a worst-slice coverage of 0.87 is not evidence of conditional miscoverage}. We report all worst-slice metrics next to this null.

\begin{proposition}[Dirichlet scale identity]\label{prop:scale}
If $y\mid x\sim\mathrm{Dirichlet}(\kappa m)$ with $m\in\Delta^{K-1}$, $\kappa>0$, then $\mathbb E\|y-m\|_2^2=\dfrac{1-\|m\|_2^2}{\kappa+1}$.
\end{proposition}
\emph{Proof.} $\mathrm{Var}(y_k)=m_k(1-m_k)/(\kappa+1)$ and $\sum_k m_k(1-m_k)=1-\|m\|^2$. \qed\\
Hence the $\ell_2$-score with oracle prediction has conditional scale $\sigma(x)\propto\sqrt{(1-\|m(x)\|^2)/(\kappa(x)+1)}$: \emph{even with constant concentration}, the scale of the score varies with the confidence $\|m(x)\|$ of the prediction (low confidence $\Rightarrow$ larger errors), and it varies with $\kappa(x)$ when concentration is heteroscedastic. For a score of the form $s=\sigma(x)\varepsilon$ with $\varepsilon\sim F$ independent of $x$, the plain threshold $q=$ quantile of $\sigma\varepsilon$ gives conditional coverage $F(q/\sigma(x))$, which is $<1-\alpha$ exactly where $\sigma(x)$ is above its $(1-\alpha)$-weighted level; normalising by $\sigma(x)$ removes the effect. Prop.~\ref{prop:scale} says the \emph{kind} of heteroscedasticity: confidence-driven (visible to the predictor through $\hat p$, hence addressable by Mondrian bins on $\|\hat p\|$) versus $\kappa$-driven (invisible in $\hat p$, addressable only by a scale model using covariates).
\paragraph{Mass and volume of prediction sets.} A set $C(x)\subset\Delta^{K-1}$ has no canonical size. We use two proxies. \emph{Mass}: the average over test $i$ and over a held-out reference sample of targets $y_j$ (drawn from the data law, independent of the test point's own target) of $\mathbf 1\{s(y_j,\hat p(x_i))/\hat\sigma_i\le\hat q_i\}$, i.e.\ the fraction of ``typical targets'' a set would accept; and for $K=3$, \emph{volume}: the same fraction with $y_j$ uniform on the simplex (4{,}000 points, first 200 test points). Mass is data-adapted; volume is measure-agnostic; we report both because rankings need not agree.

\section{Experimental Design}\label{sec:design}
\paragraph{Synthetic benchmark.} $x\sim N(0,I_5)$, $m(x)=\mathrm{softmax}(B^\top x)$ ($B$ fixed per $K$), $y\mid x\sim\mathrm{Dirichlet}(\kappa(x)m(x))$ with $\kappa(x)=\kappa_0\exp(h\,x_1)$. Factors: $K\in\{3,10\}$; $\kappa_0\in\{3,30,300\}$ (from very diffuse to concentrated); heteroscedasticity $h\in\{0,1\}$; \emph{mis-specification}: logits $B^\top x$ vs.\ $B^\top x+g(x)C$ with a nonlinear $g$, for the predictor, a soft-label multinomial logistic regression (ridge $10^{-3}$, L-BFGS) that is linear in $x$. 24 scenarios $\times$ 100 independent repetitions (new data and splits each); per repetition $n_{\rm train}=1000$, $n_{\rm cal}=500$, $n_{\rm ref}=500$, $n_{\rm test}=1000$; $\alpha=0.1$. Six scores $\times$ three calibration variants are applied to the same draw. Evaluation slices: five quantile bins of $x_1$ (the driver of $\kappa$; unobserved by the predictor's output) and of the predicted maximum probability.
\paragraph{Theory check.} For $K=3$, $\kappa_0=30$, $h=0$, correct specification, a single trained predictor is fixed; for $n\in\{20,50,100,500,2000\}$ and three scores we draw 3,000 independent calibration sets, compute the realised coverage $C$ against a $10^6$-point test pool, and compare the distribution with Prop.~\ref{prop:beta} (KS test, quantiles, $\mathbb P(C<0.9)$).
\paragraph{Real data.} CIFAR-10 test images (10,000) with CIFAR-10H label distributions \citep{peterson2019human} (about 50 annotators per image; counts in $[47,63]$; 80.6\% of entries are exactly zero, so the boundary matters). The predictor is a small 4-layer CNN trained on the 50,000 CIFAR-10 training images \emph{with hard labels only} (2 seeds, 10 epochs, one-cycle SGD; no use of CIFAR-10H), temperature-scaled on 1,000 held-out test images against the soft labels. The remaining 9,000 images supply calibration ($n_{\rm cal}\in\{100,500,2000\}$), evaluation (up to 3,000) and reference (500) sets; 30 random splits per CNN seed. Settings: \emph{exchangeable}; \emph{soft shift}: calibration images are over-sampled from confident images with the known density ratio $w(x)=e^{-3u(x)}$, $u$ the confidence rank (so weighted conformal applies); \emph{hard shift}: calibration only from the more confident half, test from the less confident half (disjoint supports; violates the assumptions of weighted conformal, so no weighted variant is run). Normaliser features: max probability, entropy, top-2 margin.
\paragraph{Uncertainty quantification and tests.} Synthetic results average over $100$ repetitions per scenario; reported standard errors treat repetitions as independent and scenarios as fixed effects, i.e.\ the SE of a grand mean over $S$ scenarios is $\frac1S\sqrt{\sum_s \mathrm{Var}_s/R}$, with scores/methods sharing a draw averaged first. For paired comparisons across methods we report scenario-level $t$-tests (the unit of replication for claims about ``across scenarios'', $n=6$ per cell) with Holm correction over the 16 comparisons, and rep-level SEs. CIFAR intervals are $\pm1.96$ split-level SEs pooled over the 2 CNN seeds; with only two CNN seeds, between-predictor variability is \emph{not} estimated beyond reporting both seeds.
